\section*{Заключение}
\addcontentsline{toc}{section}{Заключение}

В результате проведённой работы подтверждена гипотеза о том, что
отдельное веб-приложение для каталога языковых клубов и записи на
разговорные встречи решает проблему, с которой не справляются
существующие универсальные платформы мероприятий и приложения языкового
обмена. Анализ показал, что ни одно из распространённых решений
(HelloTalk, Tandem, Meetup) не сочетает одновременно групповой формат
встреч, языковую специализацию и фильтрацию по уровню владения языком~(CEFR),
что и обосновало необходимость разработки.

На основе полученных при разработке данных можно сформулировать следующие
выводы:
\begin{enumerate}
    \item Связка fullstack-фреймворка Next.js~15 и облачной платформы
          Supabase позволяет вести серверную и клиентскую части в едином
          репозитории, а аутентификацию, базу данных и файловое хранилище
          получать от одного поставщика, что заметно сокращает объём
           кода без потери гибкости.
    \item Защита данных средствами Row Level Security на уровне PostgreSQL
          является более надёжным механизмом разграничения доступа, чем
          проверки в коде сервера: правила применяются ко всем запросам
          независимо от точки входа, что исключает класс ошибок, связанных
          с пропущенными проверками прав.
    \item Хранение счётчика занятых мест прямо в таблице встреч и
          проверка лимита средствами базы данных позволяют корректно
          обрабатывать запись на встречу даже тогда, когда несколько
          пользователей нажимают кнопку <<Записаться>> одновременно:
          превышение количества свободных мест становится невозможным.
\end{enumerate}

\textbf{Рекомендации и предложения.} Для практического использования
разработанной платформы в учебных заведениях и языковых центрах предлагается
размещение на бесплатном тарифе Vercel и Supabase, что позволяет получить
функционирующий продукт без затрат на инфраструктуру. Для роста активности
рекомендуется поэтапное наполнение каталога организаторами клубов и
постепенное расширение модерации через административную панель.

\textbf{Перспективы углубления темы.} В дальнейшем развитии системы
планируется внедрение функциональности уведомлений участников о
приближающихся встречах по электронной почте и в Telegram, добавление
оценок и отзывов на проведённые встречи, поддержка иных языков помимо
английского, а также реализация мобильного клиента на базе общего API
с веб-приложением.

